\documentclass[12pt,a4paper,oneside]{article}
\usepackage[brazil]{babel}
\usepackage[utf8]{inputenc}
\usepackage{olymp}

\setcounter{problem}{1}

\begin{document}

\begin{problem}{Doublets}{doublets}{2}
 
Um \textit{doublet} é formado por um par de palavras de mesmo tamanho que diferem entre si em somente uma letra. Por exemplo: ``bola'' e ``bala'', ou ``passa'' e ``massa''. Faça um programa para determinar se duas cadeias de caracteres formam um \textit{doublet} ou não.

%\Notes
%
%Observações, se tiver.

\InputFile

A entrada contém exatamente duas linhas, cada uma delas com uma cadeia de no máximo 100 caracteres, contendo somente letras minúsculas.

\OutputFile

Seu programa deve escrever somente ``\texttt{SIM}'' ou ``\texttt{NAO}' (sem acento), indicando se as duas cadeias formam ou não um \textit{doublet}.
% \Notes
% Preste atenção aos tipos das variáveis, pois $F(90) \approx 2^{61}$.

\Example

\begin{example}
\exmp{
bolo
bola
}{
SIM
}%
\end{example}

\begin{example}
\exmp{
passa
massa
}{
SIM
}%
\end{example}

\begin{example}
\exmp{
marra
massa
}{
NAO
}%
\end{example}

%Comentários sobre o exemplo, se necessário.

\end{problem}

\end{document}
